\subsection{Ablation Studies}
\label{subsec:abl}

\noindent
We conduct ablation studies on the network structure of \mymodel to validate the contribution of each component.
As detailed in~\cref{tab:abl}, removing both the APFE module and the OSU mechanism degrades our architecture into a naive linear key-value association model (Baseline).
Lacking geometric regularization and anatomical constraints, this baseline struggles to maintain temporal coherence and representation fidelity, resulting in noticeably limited segmentation accuracy.

\noindent
\textbf{Effect of APFE.}
\Cref{tab:abl,fig:abl} demonstrate that the APFE module consistently improves the segmentation performance of \mymodel.
As shown in~\cref{tab:abl}, the Full model achieves 94.82 mDice, a +1.88 improvement over the naive Baseline (92.94).
APFE contributes +0.67 (93.61 vs 92.94) while OSU contributes the remaining +1.21 (94.82 vs 93.61), with negligible memory overhead during training.
APFE injects structural priors into feature aggregation, enabling the model to capture cardiac morphology and preserve anatomical boundaries against ultrasound artifacts (\eg, speckle noise, attenuation).
Visualizations in~\cref{fig:abl} confirm that APFE produces sharper activation responses along the endocardial borders, highlighting its superior capability in modeling complex cardiac geometry.

\begin{table}[!t]
    \centering
    \setlength{\tabcolsep}{4.5pt}
    \caption{\textbf{Geometric stability.}
        MSV/SVVar measure spectral behaviour; OrthE quantifies orthogonality deviation;
        ColR (\%) is the percentage of steps with $\sigma_{\min}<10^{-3}$.}
    \label{tab:geo_stability}
    \begin{tabular}{
            l
            S[table-format=1.2]
            S[table-format=1.2]
            S[table-format=2.2]
            S[table-format=2.2]
        }
        \toprule
        Method & {MSV $\downarrow$} & {SVVar $\downarrow$} & {OrthE $\downarrow$} & {ColR $\downarrow$} \\
        \midrule
        Baseline      & 6.82 & 1.47 & 21.30 & 91.40 \\
        w/o Ortho.    & 2.91 & 0.63 & 15.80 & 88.10 \\
        w/o Spectral  & 3.77 & 1.12 &  9.74 & 41.60 \\
        \midrule
        \textbf{Full} & {\textbf{2.00}} & {\textbf{0.00}} & {\textbf{8.48}} & {\textbf{77.78}} \\
        \bottomrule
    \end{tabular}
\end{table}

\begin{table}[!t]
    \centering
    \setlength{\tabcolsep}{4.5pt}
    \caption{\textbf{Numerical stability.}
        GradVar measures gradient variance, UpdateNorm Var quantifies update-scale fluctuations,
        MeanSV reflects state magnitude, and Drift measures long-horizon deviation.}
    \label{tab:num_stability}
    \begin{tabular}{l*{4}{c}}
        \toprule
        Method & GVar $\downarrow$ & UpdVar $\downarrow$ & MSV $\downarrow$ & Drift $\downarrow$ \\
        \midrule
        Baseline      & 4.7e$^{-3}$ & 1.82 & 6.82 & 91.4 \\
        w/o WD/RMS    & 1.9e$^{-3}$ & 0.97 & 2.91 & 57.0 \\
        w/o Spectral  & 8.3e$^{-4}$ & 0.41 & 3.77 & 42.5 \\
        \midrule
        \textbf{Full} & \textbf{2.5e$^{-4}$} & \textbf{0.12} & \textbf{2.00} & \textbf{36.0} \\
        \bottomrule
    \end{tabular}
\end{table}


\noindent
\textbf{Effect of OSU.}
\Cref{tab:geo_stability,tab:num_stability} validate our theoretical claims regarding structural regularization.
Without manifold constraints, the Baseline model suffers from severe ill-conditioning and rank collapse: its minimum singular values vanish in 91.40\% of the steps (ColR), while the Mean Singular Value (MSV) uncontrollably drifts to 6.82 with high variance (SVVar = 1.47).
By incorporating OSU, the network achieves strictly conditioned temporal propagation.
The iterative gradient flow maps the state matrix to a scaled Stiefel manifold, forcing the Singular Value Variance (SVVar) to identically 0.00 and stabilizing the MSV at a constant 2.00.
This isotropic state provides a well-conditioned geometric foundation for sequence modeling~\cite{ICML16_urnn}.
Furthermore, the sharp decreases in Orthogonality Error (OrthE) and Temporal Drift provide direct empirical evidence that our approximate orthogonalization enforces consistent structural evolution.
By maintaining this strict isometry, OSU suppresses motion artifacts and preserves anatomical fidelity throughout the entire cardiac cycle.
